\section{状态空间、确定性与噪声}

\begin{definition}[状态空间描述]
\textbf{状态空间描述}用一个状态向量 $x(t)$ 加上一条更新规则来表示系统。
在连续时间中，可写为
\[
  \dot{x} = f(x,u,\eta),
\]
其中 $x$ 表示当前状态，$u$ 表示控制或干预输入，$\eta$ 表示随机性影响或未建模影响。
\end{definition}

\begin{remark}
这个方程并不是作为一个已经拟合好的人类模型被提出的。
它是一个有纪律的模板。
它告诉读者哪些变量类型重要：当前状态、外部输入与噪声。
仅此一点，就已经比模糊地说“我感觉卡住了”前进了一大步。
\end{remark}

这一记号还有助于整合前面几章。
第~\ref{ch:consciousness-biology} 章为在生物学中讨论有组织的状态转变提供了强理由。
第~\ref{ch:emergence} 章说明，重复的局部互动能够形成一个对历史敏感的机制。
本章则补上必需的数学语言，用来说明那种机制可能属于何种类型：
一个不动点、一个亚稳态、一个极限环、一条带噪声的转变路径，或者系统行为中一条对控制输入敏感的分支。